% Curriculum vitae, carrying the same material as cv.html on the website.
% Build with tools/build_cv.sh, which writes files/CV_Miguel_Zumalacarregui.pdf.
% This directory is underscore-prefixed so Jekyll leaves the source unpublished.
%
% Keep this in step with cv.html: the page is the source the site renders, and
% this file is what people download. Sections appear here in the page's order.
\documentclass[11pt,a4paper]{article}

\usepackage[T1]{fontenc}
\usepackage[utf8]{inputenc}
\usepackage[margin=2.2cm,top=1.9cm,bottom=1.9cm]{geometry}
\usepackage{enumitem}
\usepackage{titlesec}
\usepackage{xcolor}
\usepackage[hidelinks]{hyperref}

\definecolor{accent}{HTML}{4B3FA8}
\definecolor{soft}{HTML}{444444}

\hypersetup{
  colorlinks=true,
  urlcolor=accent,
  pdftitle={Curriculum vitae — Miguel Zumalacárregui},
  pdfauthor={Miguel Zumalacárregui},
  pdfsubject={Curriculum vitae}
}

\pagestyle{plain}
\setlength{\parindent}{0pt}

% Section headings: small caps over a rule, matching the site's quiet hierarchy.
\titleformat{\section}
  {\normalfont\large\bfseries\color{accent}}{}{0pt}{}[\vspace{-0.6em}\rule{\textwidth}{0.6pt}]
\titlespacing*{\section}{0pt}{1.15em}{0.55em}

\titleformat{\subsection}{\normalfont\normalsize\bfseries}{}{0pt}{}
\titlespacing*{\subsection}{0pt}{1.1em}{0.35em}

% One entry: bold headline, then a quieter line of institution, place and dates.
\newcommand{\entry}[2]{%
  \item {#1}%
  \ifx\relax#2\relax\else\\[0.1em]{\small\color{soft}#2}\fi
}
\newenvironment{entries}
  {\begin{itemize}[leftmargin=1.1em,itemsep=0.32em,parsep=0pt,topsep=0.3em]}
  {\end{itemize}}
\newenvironment{plainlist}
  {\begin{itemize}[leftmargin=1.1em,itemsep=0.2em,parsep=0pt,topsep=0.3em]}
  {\end{itemize}}

\begin{document}

% ---------------------------------------------------------------- header ----
\begin{center}
  {\LARGE\bfseries Miguel Zumalacárregui}\\[0.45em]
  {\color{soft}Group Leader, Astrophysical and Cosmological Relativity ---
   Max Planck Institute for Gravitational Physics (Albert Einstein Institute), Potsdam}\\[0.2em]
  {\color{soft}Ramón y Cajal Fellow --- Instituto de Física Teórica IFT-UAM/CSIC, Madrid}\\[0.5em]
  {\small
   \href{mailto:miguel.zumalacarregui@aei.mpg.de}{miguel.zumalacarregui@aei.mpg.de}
   \quad$\cdot$\quad
   \href{https://miguelzuma.github.io}{miguelzuma.github.io}
   \quad$\cdot$\quad
   \href{https://orcid.org/0000-0002-9943-6490}{ORCID 0000-0002-9943-6490}}
\end{center}

\vspace{0.2em}
{\small\color{soft}
The publication record is maintained on
\href{https://inspirehep.net/authors/1113379}{INSPIRE-HEP}; research
\href{https://miguelzuma.github.io/projects.html}{projects and funding} and the
\href{https://miguelzuma.github.io/group.html}{people supervised} are listed on the website.}

% ------------------------------------------------------------- positions ----
\section{Positions}
\begin{entries}
  \entry{\textbf{Ramón y Cajal Fellow}}
        {Instituto de Física Teórica IFT-UAM/CSIC, Madrid $\cdot$ 2026--}
  \entry{\textbf{Group Leader}, Astrophysical and Cosmological Relativity}
        {Max Planck Institute for Gravitational Physics (Albert Einstein Institute), Potsdam $\cdot$ 2020--}
  \entry{\textbf{Marie Skłodowska-Curie Global Fellow}}
        {Berkeley Center for Cosmological Physics, with return phase at IPhT Saclay $\cdot$ 2017--2020}
  \entry{\textbf{Nordita Fellow}}
        {Nordic Institute for Theoretical Physics, Stockholm $\cdot$ 2015--2017}
  \entry{\textbf{Postdoctoral Researcher}}
        {Institut für Theoretische Physik, University of Heidelberg $\cdot$ 2013--2015}
  \entry{\textbf{Postdoctoral Researcher}}
        {Instituto de Física Teórica IFT-UAM/CSIC, Madrid $\cdot$ 2013}
\end{entries}

% ------------------------------------------------------------- education ----
\section{Education}
\begin{entries}
  \entry{\textbf{PhD in Physics}, \emph{Cum Laude} with International Mention}
        {Universidad Autónoma de Madrid $\cdot$ 2012 $\cdot$ Graduate work at ICC Barcelona and IFT-UAM/CSIC (FPI fellowship)\\
         Dissertation: \href{https://miguelzuma.github.io/files/PhD_thesis_miguel_zuma.pdf}{Probing the Foundations of the Standard Cosmological Model}}
  \entry{\textbf{Master's thesis}}
        {Universidad Autónoma de Madrid $\cdot$
         \href{https://miguelzuma.github.io/files/master_thesis_miguel_zuma.pdf}{PDF}}
\end{entries}

% ----------------------------------------------------------- supervision ----
% Team members only: the `current` and `alumni` lists of _data/people.yml.
% Visiting researchers and the informal mentees from before the Potsdam group
% are left to the website, which is where the full record lives.
\section{Supervision}

\subsection{Postdoctoral researchers}
\begin{plainlist}
  \item \textbf{Han-Gil Choi} --- 2026--
  \item \textbf{Xikai Shan} --- affiliated postdoctoral researcher, 2026--
  \item \textbf{Juan Urrutia} --- from October 2026
  \item \textbf{Héctor Villarrubia-Rojo} --- from October 2026; previously Margarita Salas Fellow in the group, 2022--2024
  \item \textbf{Daniel Ballard} --- from October 2026
  \item \textbf{Srashti Goyal} --- 2023--2026 $\rightarrow$ Independent Fellow at IUCAA
  \item \textbf{Serena Giardino} --- 2024--2025 $\rightarrow$ postdoc at Imperial College London
  \item \textbf{Guilherme Brando} --- Humboldt Fellow, 2022--2024 $\rightarrow$ tenured researcher at UFES, Brazil
  \item \textbf{Giovanni Tambalo} --- 2020--2023 $\rightarrow$ postdoc at ETH Zürich, then Marie Skłodowska-Curie Fellow at IPhT Saclay
\end{plainlist}

\subsection{PhD students}
\begin{plainlist}
  \item \textbf{Manon van Zyl} --- from October 2026
  \item \textbf{Emanuela Alberici} --- from October 2026
  \item \textbf{Nicola Menadeo} --- co-supervised, U. Naples, 2023--2025 $\rightarrow$ postdoc at Humboldt University of Berlin
  \item \textbf{Stefano Savastano} --- co-supervised with A. Buonanno, 2020--2024 $\rightarrow$ PhD awarded at HU Berlin
\end{plainlist}

\subsection{Master's and undergraduate students}
\textbf{Pablo Vega del Castillo} (LMU, 2026); \textbf{Matteo Pagani} (co-supervised, U. Bologna, 2026);
\textbf{Louis Salimbeni} (co-supervised, ETH Zürich, 2026); \textbf{Selin Üstündag} (2026, now PhD at
Newcastle University); \textbf{Marcel Glätzer} (HU Berlin, 2024--2025); \textbf{Jaime Raynaud}
(TU Berlin \& KTH, 2024--2025, now PhD at MPIA); \textbf{Daryna Yushchenko} (Princeton summer
programme, 2024); \textbf{Vaibhav Omble} (LMU, 2023--2024); \textbf{Julius Streibert} (co-supervised,
FU Berlin, 2022--2023); \textbf{Shashwat Singh} (Paris Observatory, 2022--2023, now PhD at the
University of Glasgow); \textbf{Lyla Choi} (Princeton summer programme, 2023, now PhD at Stanford);
\textbf{Richard Stiskalek} (co-supervised, LMU, 2021--2022, now postdoc at Oxford);
\textbf{Alessandro Longo} (co-supervised, U. Trieste, 2020--2021, now PhD student at SISSA).

\vspace{0.4em}
{\small\color{soft}Visiting researchers, and the students mentored informally before the group
started in 2020, are listed at
\href{https://miguelzuma.github.io/group.html}{miguelzuma.github.io/group.html}.}

% ---------------------------------------------------------- coordination ----
\section{Scientific coordination}
\begin{entries}
  \entry{\textbf{LISA Consortium} --- member since 2018}
        {Council member and Senior Co-Chair of the Cosmology Working Group (2025--); section
         coordinator for gravitational-wave lensing in the LISA Cosmology white paper (2022)}
  \entry{\textbf{Euclid Consortium} --- member since 2014}
        {Lead of the Science Working Group on Gravitational Waves and of its ``bright sirens''
         work package (2025--); member of the Strong Lensing, Cosmology Theory, and Supernovae
         \& Transients working groups}
  \entry{\textbf{GW-Space 2050}}
        {Coordinator for the cosmography section, defining ESA's gravitational-wave missions
         after LISA (2023--)}
  \entry{\textbf{Multi-messenger lensing white paper}}
        {Coordinator for the microlensing section (2024)}
\end{entries}

% ------------------------------------------------------------- workshops ----
\section{Workshops and meetings organized}
\begin{entries}
  \entry{\href{https://indico.icc.ub.edu/event/677/}{13th LISA Cosmology Working Group Workshop}}
        {Barcelona, 2026}
  \entry{\href{https://indico.cern.ch/event/1488828/}{12th LISA Cosmology Working Group Workshop}}
        {Tallinn, 2025}
  \entry{\href{https://www.esi.ac.at/events/e546/}{Lensing and Wave Optics in Strong Gravity}}
        {Erwin Schrödinger Institute, Vienna, 2024}
  \entry{\href{https://workshops.aei.mpg.de/fpmeetswavelisa/}{LISA joint Fundamental Physics + Waveform meeting}}
        {Potsdam, 2024}
  \entry{\href{https://www.aei.mpg.de/139117/seminars}{ACR Division Seminars}}
        {Max Planck Institute for Gravitational Physics, 2020--2022}
  \entry{BCCP Cosmology Lunch}{UC Berkeley, 2016--2017}
  \entry{Cosmology Seminars, Institut für Theoretische Physik}{University of Heidelberg, 2014--2015}
  \entry{TRR33 Winter Schools}{Passo del Tonale, 2013, 2014 and 2015}
\end{entries}

% -------------------------------------------------------------- teaching ----
\section{Teaching}

\subsection{Graduate schools}
\begin{entries}
  \entry{\textbf{Gravitational waves} --- TAE (Taller de Altas Energías), Benasque, 2024}
        {\href{https://www.benasque.org/2024tae/talks_contr/094_GW_Lectures_-_TAE_2024_-_Zumalacarregui_-_PDF.pdf}{Slides}}
  \entry{\textbf{Tests of gravity and dark energy with cosmology and gravitational waves}}
        {II Mesoamerican Workshop on Gravity and Cosmology, Chiapas, 2018}
  \entry{\textbf{Alternative gravity theories} --- Cosmology School in the Canary Islands, 2017}
        {\href{http://iactalks.iac.es/talks/view/1081}{Video lecture}}
  \entry{\textbf{School for Cosmology Tools} --- IFT Madrid, 2017}
        {Lectures on CLASS, hi\_class and MontePython}
  \entry{\textbf{Workshop on CLASS and MontePython} --- MPIA Munich, 2014}{}
\end{entries}

\subsection{University courses}
\begin{entries}
  \entry{\textbf{Berkeley Connect Instructor} --- UC Berkeley, 2018--2019}
        {Two semesters of undergraduate mentoring and teaching}
  \entry{\textbf{Lecturer} --- University of Heidelberg, 2014--2015}
        {Advanced Cosmology; Introduction to CMB Physics; Statistical Mechanics; advanced seminar}
  \entry{\textbf{Teaching assistant} --- University of Heidelberg, 2014--2015}
        {Theoretical Astrophysics (graduate), with M. Bartelmann}
  \entry{\textbf{Teaching assistant} --- Universidad Autónoma de Madrid, 2013}
        {Gravitation and Cosmology; Physics of the Cosmos}
  \entry{\textbf{Mini-course on inhomogeneous cosmologies} --- University of Barcelona, 2012}{}
\end{entries}

% ---------------------------------------------------------------- awards ----
\section{Awards and fellowships}
\begin{plainlist}
  \item Ramón y Cajal Fellowship (2024).
  \item R3 certification, Spanish Research Agency.
  \item Berkeley Connect Fellow, UC Berkeley (2018).
  \item Marie Skłodowska-Curie Global Fellowship (2015).
  \item Nordita Fellowship (2015).
  \item Yggdrasil Fellowship (2011) --- research stay at ITA Oslo.
  \item FPI doctoral fellowship (2008) and UAM Master's scholarship (2007).
\end{plainlist}

% ----------------------------------------------------------- peer review ----
\section{Evaluation and peer review}

\subsection{Evaluation panels}
\begin{plainlist}
  \item NASA Hubble Fellowship Program, Physics and Cosmology panel.
  \item FCT --- Portuguese Foundation for Science and Technology, Individual Support Call.
  \item Max Planck Institute for Gravitational Physics --- postdoctoral hiring committee, since 2020.
\end{plainlist}

\subsection{External evaluator for funding agencies}
European Research Council (Starting Grants), STFC/UKRI, Swiss National Science Foundation,
Alexander von Humboldt Foundation, InterTalentum (Marie Skłodowska-Curie COFUND),
Netherlands Research Council (NWO), Research Grants Council of Hong Kong, and the
Estonian Research Council.

\subsection{Refereeing}
Peer review for more than 40 papers, including Physical Review Letters, Physical Review D,
Nature Astronomy, JCAP, The Astrophysical Journal and Astrophysical Journal Letters.

% ----------------------------------------------------------------- talks ----
\section{Selected talks}
\begin{entries}
  \entry{\textbf{Spanish \& Portuguese Relativity Meeting} --- Universidad de Murcia, 2026 $\cdot$ plenary}
        {Gravitational lensing of waves: a new window into astrophysics, dark matter and gravity}
  \entry{\textbf{IAS and Princeton University} --- gravitational-wave meeting, Princeton, 2026 $\cdot$ seminar}
        {GW231123 as a diffracted and magnified black-hole merger}
  \entry{\textbf{Leibniz Institute for Astrophysics} --- Potsdam, 2024 $\cdot$ institute colloquium}
        {Gravitational lensing of waves}
  \entry{\textbf{New Frontiers in Strong Gravity} --- Benasque, 2024 $\cdot$ plenary}
        {Wave-optical phenomena in gravitational-wave lensing}
  \entry{\textbf{Workshop on Tensions in Cosmology} --- Corfu, 2022 $\cdot$ plenary}
        {Towards solutions to the Hubble problem beyond Einstein's gravity}
  \entry{\textbf{EuCAPT Annual Symposium} --- CERN, 2022 $\cdot$ plenary review}
        {\href{https://cds.cern.ch/record/2810852}{Dark energy and tests of gravity with cosmology and gravitational-wave propagation}}
  \entry{\textbf{Institute of Theoretical Astrophysics} --- University of Oslo, 2021 $\cdot$ institute colloquium}
        {The dark Universe in the era of gravitational-wave astronomy}
  \entry{\textbf{Progress in Old and New Themes in Cosmology} --- Palais des Papes, Avignon, 2020 $\cdot$ plenary}
        {Tests of gravity with gravitational waves}
  \entry{\textbf{COSMO} --- Aachen University, 2019 $\cdot$ plenary}
        {Tests of dark energy and modified gravity}
  \entry{\textbf{Rencontres du Vietnam} --- Very High Energy Phenomena in the Universe, 2018 $\cdot$ plenary}
        {The dark Universe in the era of gravitational-wave astronomy}
\end{entries}

% ------------------------------------------------------------- languages ----
\section{Languages}
Spanish (native), English (C2), German (C1).

\end{document}
